\section{Problema productor-consumidor}

El \textit{problema productor-consumidor} modela la comunicación entre procesos mediante un \textit{búfer compartido} de tamaño finito. Uno o más procesos \textit{productores} generan elementos y los insertan en el búfer, mientras que uno o más procesos \textit{consumidores} los extraen para procesarlos.

La sincronización debe garantizar que:

\begin{itemize}
  \item Un productor no inserte elementos cuando el búfer esté lleno.
  \item Un consumidor no extraiga elementos cuando el búfer esté vacío.
  \item El acceso al búfer sea mutuamente excluyente.
\end{itemize}

Para resolver el problema se utilizan tres semáforos:

\begin{itemize}
  \item \texttt{mutex = 1}: protege el acceso al búfer (exclusión mutua).
  \item \texttt{empty = N}: indica la cantidad de posiciones libres del búfer.
  \item \texttt{full = 0}: indica la cantidad de posiciones ocupadas.
\end{itemize}

El productor espera a que exista una posición libre, accede al búfer de forma exclusiva para insertar un elemento y, posteriormente, señala que existe un nuevo elemento disponible.

\begin{ccode}[title=Implementación del productor]
semaphore mutex = 1;
semaphore empty = N;
semaphore full = 0;

void producer() {
  while (true) {
    item = produce();

    semWait(empty); // decrementa empty
    semWait(mutex); // decrementa mutex

    insert(item);

    semSignal(mutex); // incrementa mutex
    semSignal(full); // incrementa full
  }
}
\end{ccode}

El consumidor espera hasta que exista al menos un elemento disponible, accede al búfer de forma exclusiva para retirarlo y, finalmente, señala que una posición ha quedado libre.

\begin{ccode}[title=Implementación del consumidor]
void consumer() {
  while (true) {
    semWait(full); // decrementa full
    semWait(mutex); // decrementa mutex

    item = remove();

    semSignal(mutex); // incrementa mutex
    semSignal(empty); // incrementa empty

    consume(item);
  }
}
\end{ccode}

Este algoritmo garantiza que ningún productor escriba sobre un búfer lleno, ningún consumidor lea de un búfer vacío y que nunca dos procesos accedan simultáneamente al búfer compartido.
